\documentclass[a4paper]{article}

\usepackage{graphicx}
\usepackage{amsmath,amssymb}
\usepackage{minted}
\usepackage{multirow}
\usepackage{float}
\usepackage{algorithm}
\usepackage{algpseudocode}
\usepackage{todonotes}
\usepackage{forest}
\usepackage{xstring}
\usepackage{xcolor}
\usepackage[most]{tcolorbox}
\usepackage{xparse}
\usepackage{natbib}

\title{Selection-only bookmark range fix}
\date{2026-09-18}
\author{Mohammad Rahmani}

\begin{document}
\maketitle

\section{Problem}
A short text selection could resolve through the Visual line/column approximation
to a much broader canonical range, so a bookmark could cover a paragraph instead
of only the selected text. Legacy cursor-oriented entry points also made the
bookmark semantics look line-based even though the persisted model was range-based.

The persistent Visual bookmark paint also depended on the same approximate
line/column projection. Headings and text following inline LaTeX could therefore
lose the exact selected highlight after the normal QTextEdit selection changed.

\section{Fix}
Bookmark creation is now strictly selection-only. Source selections use their
exact character offsets, including the canonical offset of a virtualized Source
window. Visual selections first resolve the exact selected LaTeX or visible text
near the projected position; an implausibly broad fallback is rejected rather
than converted into a paragraph, line, or structural bookmark.

Bookmark toggle semantics now compare exact ranges. Merely overlapping an
existing bookmark no longer means that the existing bookmark is the selected
bookmark. The compatibility cursor toggle delegates to selection-only behavior,
so a cursor without a selection cannot create a bookmark.

A selection-aware bookmark highlighter now resolves the stored bookmarked text
back into the Visual document near its canonical projection. The yellow
highlight therefore remains on the exact bookmarked text after the transient
selection is released, including ordinary headings and inline-formatted text.

\section{Regression coverage}
Tests cover exact regular Source offsets, virtualized Source offsets, Visual
selection recovery from an incorrectly broad paragraph projection, rejection of
unsafe broad fallbacks, exact-range toggle semantics, the no-selection case, and
exact Visual bookmark text lookup.

\end{document}
